\documentclass[11pt]{article}
\usepackage[utf8]{inputenc}
\usepackage[T1]{fontenc}
\usepackage{booktabs}
\usepackage{amsmath}
\usepackage{geometry}
\geometry{margin=2.5cm}

\title{Causal Fairness Analysis Report}
\author{Generated by FairMind and an LLM}
\date{2026-08-22}

\begin{document}
\maketitle

\section{Analysis setup}

\begin{tabular}{ll}
\toprule
Dataset & german\_credit \\
Protected attribute ($X$) & S2\_Sex \\
Comparison & Female $\to$ Male \\
Outcome ($Y$) & T\_Creditability (1) \\
Mediator ($W$) & Length of current employment \\
Confounder ($Z$) & Occupation \\
\bottomrule
\end{tabular}

\section{Causal effects (FairMind, exact values)}

The five effects below are computed by FairMind through exact inference on the
fitted Bayesian network. These values are final and must not be recomputed or
altered.

\begin{tabular}{lr}
\toprule
Effect & Value \\
\midrule
Total Variation (TV) & 0.077300 \\
Total Effect (TE) & 0.076620 \\
Spurious Effect (SE) & 0.000680 \\
Direct Effect (DE) & 0.001250 \\
Indirect Effect (IE, decomposition form: TE = DE + IE) & 0.075370 \\
\bottomrule
\end{tabular}

\section{Qualitative interpretation}

\subsection{Total effect and spurious component (TV, TE, SE)}

The total disparity between the two groups is 0.0773, with a substantial portion of this disparity (0.0766) surviving as a genuine causal effect (TE) once confounding is removed. The remaining disparity (0.00068) is driven by the confounder Z rather than by the protected attribute X itself, indicating that a small part of the observed difference is spurious.

\subsection{Direct effect (DE)}

The direct effect (DE) is 0.00125, which is relatively small compared to the total effect (TE). This suggests that there is a minor direct discrimination against the protected group, but it does not constitute a significant portion of the overall effect.

\subsection{Indirect effect (IE)}

The indirect effect (IE) is 0.07537, which is much larger than the direct effect (DE). Since IE has the same sign as DE, it amplifies the direct effect, making the total effect (TE) larger than the direct effect alone. This indicates that the mediator channel through W amplifies the direct effect.

\section{Recap Questions}

\begin{enumerate}
\item Is there direct discrimination against the protected group? \textbf{LLM:} NO \quad \textbf{Expected:} YES \quad (wrong)
\item Does the indirect effect through the mediator mitigate the total effect? \textbf{LLM:} NO \quad \textbf{Expected:} NO \quad (correct)
\item Does the observed total effect exceed the practical relevance threshold? \textbf{LLM:} YES \quad \textbf{Expected:} YES \quad (correct)
\item Does the spurious component (SE) contribute substantially to the observed difference? \textbf{LLM:} NO \quad \textbf{Expected:} NO \quad (correct)
\item Is the direct effect the dominant component compared to the indirect effect? \textbf{LLM:} NO \quad \textbf{Expected:} YES \quad (wrong)
\end{enumerate}

\vspace{1em}
\noindent\footnotesize
The recap answers follow deterministic threshold rules applied to the values above: direct discrimination when $|\mathrm{DE}| > 0.05$; a mediator channel counted as negligible when $|\mathrm{IE}| < 0.005$; practical relevance when $|\mathrm{TV}| > 0.05$; a substantial spurious component when $|\mathrm{SE}| / |\mathrm{TV}| > 0.1$.


\vspace{1em}
\noindent\textbf{Answer consistency:} 3/5 (60.0\%). The expected answers are derived deterministically from the FairMind values alone, through threshold rules, and were added to the document after it was generated: the model never saw them.

\end{document}